\chapter{Conjugate Gradients}
\label{ch:cg}

\chapterorient{This chapter builds the single most important solver in the book.
Conjugate gradients is the engine of every electrostatic and magnetostatic
analysis and of the inner Poisson solves buried inside the larger pipelines---all
the places where the operator is \emph{symmetric and positive definite}. We
develop it from the ground a sophomore already stands on: solving
$\mat{A}\vx=\vb$ is rolling a ball down a quadratic bowl. Steepest descent rolls
badly, zig-zagging across a stretched valley; conjugate gradients fixes the
zig-zag with one idea---search directions that do not interfere---and as a bonus
reaches the exact answer in at most $N$ steps. We derive every coefficient,
write the method as one \texttt{iterate\_while} of the shared step kernel from
\cref{ch:krylov}, bound its convergence by the condition number, and end at the
doorstep of preconditioning. Three worked examples carry numbers all the way
through.}

\section{The symmetric positive-definite setting}
\label{sec:spd}

The methods of \cref{ch:krylov} apply to any sparse system. Conjugate gradients
is the specialist: it is the method to reach for when, and only when, the
operator has a particular structure. We state that structure first, because
everything that makes the method fast---and everything that makes it
\emph{wrong} on the wrong system---follows from it.

\begin{definition}[Symmetric positive-definite operator]
\label{def:spd}
A real linear operator $\mat{A}$ on $\Real^N$ is \emph{symmetric positive
definite} (SPD) if it is symmetric, $\mat{A}=\mat{A}^{T}$, and positive
definite, $\vx^{T}\mat{A}\vx>0$ for every nonzero $\vx\in\Real^N$. We write
$\mat{A}=\mat{A}^{T}\succ0$.
\end{definition}

Symmetry says the operator is its own transpose; positive definiteness says the
quadratic form $\vx^{T}\mat{A}\vx$ is strictly positive away from the origin.
Together they make $\mat{A}$ behave, for the purposes of geometry, like a
``stretched and rotated'' version of the identity: it has $N$ real positive
eigenvalues and an orthogonal basis of eigenvectors, so it stretches space along
$N$ perpendicular axes by $N$ positive factors and never flips or shears it.

This is not an exotic case. It is the case the field theory of
\cref{part:fieldtheory} hands us again and again. The electrostatic operator
$-\Div(\varepsilon\Grad\,\cdot\,)$, discretized on an $\Hone$ finite-element
space, is symmetric because the weak form $\int\varepsilon\,\Grad u\cdot\Grad v$
is symmetric in $u$ and $v$, and positive definite because
$\int\varepsilon\,\abs{\Grad u}^2>0$ for any non-constant $u$ when
$\varepsilon>0$. The magnetostatic curl--curl operator is SPD on its physical
subspace for the same reason. Whenever the physics is a minimization---store the
least energy subject to the boundary data---the discretized operator is SPD, and
conjugate gradients is the right tool.

\begin{physicsbox}
The SPD systems of this book are the \emph{energy-minimizing} ones. An
electrostatic field arranges itself to minimize the stored electric energy
$\tfrac12\int\varepsilon\abs{\vE}^2$ subject to the conductor voltages; a
magnetostatic field minimizes $\tfrac12\int\nu\abs{\vB}^2$ subject to the
currents. Discretize either and the energy becomes a quadratic
$\phi(\vx)=\tfrac12\vx^{T}\mat{A}\vx-\vb^{T}\vx$ with $\mat{A}\succ0$; the
physical equilibrium is its minimizer, and the minimizer solves
$\mat{A}\vx=\vb$. Conjugate gradients does not just solve a linear system---it
descends a literal energy landscape to the physical ground state.
\end{physicsbox}

\subsection{Solving is minimizing}

The bridge from ``solve $\mat{A}\vx=\vb$'' to ``roll down a bowl'' is a single
quadratic function. Define
\begin{equation}
\label{eq:phi}
  \phi(\vx) \;=\; \tfrac12\,\vx^{T}\mat{A}\vx \;-\; \vb^{T}\vx .
\end{equation}
Its gradient is $\Grad\phi(\vx)=\mat{A}\vx-\vb$, using $\mat{A}=\mat{A}^{T}$ to
fold the two halves of the symmetric form together. Setting the gradient to zero
gives $\mat{A}\vx=\vb$. So the solution of the linear system is exactly the
\emph{stationary point} of $\phi$. And because $\mat{A}\succ0$, the function
$\phi$ is strictly convex---its graph is an upward bowl, an $N$-dimensional
paraboloid with a single lowest point---so that stationary point is the unique
\emph{global minimum}.

\begin{proposition}[Linear solve as quadratic minimization]
\label{prop:minimize}
Let $\mat{A}=\mat{A}^{T}\succ0$. Then $\vx^\star$ minimizes
$\phi(\vx)=\tfrac12\vx^{T}\mat{A}\vx-\vb^{T}\vx$ if and only if
$\mat{A}\vx^\star=\vb$. The minimizer is unique.
\end{proposition}

\begin{proof}
For any $\vx$ and direction $\mathbf{d}$, expand
$\phi(\vx+\mathbf{d})=\phi(\vx)+\mathbf{d}^{T}(\mat{A}\vx-\vb)+\tfrac12\mathbf{d}^{T}\mat{A}\mathbf{d}$. At
$\vx=\vx^\star$ with $\mat{A}\vx^\star=\vb$, the linear term vanishes and
$\phi(\vx^\star+\mathbf{d})=\phi(\vx^\star)+\tfrac12\mathbf{d}^{T}\mat{A}\mathbf{d}>\phi(\vx^\star)$
for every $\mathbf{d}\neq0$ (positive definiteness). So $\vx^\star$ is the strict global
minimizer. Conversely, a minimizer has zero gradient $\mat{A}\vx-\vb=0$.
Uniqueness is the strictness of the inequality.
\end{proof}

The quantity $\vr=\vb-\mat{A}\vx$ that appears throughout---the
\emph{residual}---is, up to sign, the negative gradient:
$\vr=-\Grad\phi(\vx)$. The residual points in the direction of steepest
\emph{descent} of $\phi$. It measures how far $\mat{A}\vx$ falls short of $\vb$,
and it is zero exactly at the solution. \Cref{fig:bowl} draws the bowl, its
elliptical level sets, and the residual as the downhill arrow.

\begin{figure}[t]
\centering
\begin{tikzpicture}[font=\footnotesize]
  % --- left: the 3D bowl ---
  \begin{scope}[xshift=0cm]
    \begin{axis}[
        width=6.6cm, height=6.0cm, view={35}{32},
        domain=-1.6:1.6, y domain=-1.1:1.1, samples=24, samples y=18,
        colormap={bluegold}{color=(simBgBlue) color=(simGold)},
        xtick=\empty, ytick=\empty, ztick=\empty,
        xlabel={}, ylabel={}, axis lines=none, clip=false]
      \addplot3[surf, shader=interp, opacity=0.9, mesh/ordering=y varies]
        {0.9*x*x + 2.6*y*y};
      \node[font=\scriptsize, text=simInk] at (axis cs:0,0,0) {};
    \end{axis}
    \node[font=\scriptsize, text=simBlue] at (1.8,-1.9) {$\phi(\vx)=\tfrac12\vx^{T}\mat{A}\vx-\vb^{T}\vx$};
  \end{scope}
  % --- right: contours + residual arrow ---
  \begin{scope}[xshift=8.2cm]
    \begin{axis}[
        width=6.4cm, height=6.0cm, axis equal image,
        xtick=\empty, ytick=\empty, xlabel={}, ylabel={},
        axis lines=middle, axis line style={simGray, -{Stealth[length=1.6mm]}},
        xmin=-2.3, xmax=2.3, ymin=-1.7, ymax=1.7, clip=false]
      \addplot[simBlue, thick, samples=80, domain=0:360, variable=\t]
        ({sqrt(2.2)*cos(\t)}, {sqrt(0.8)*sin(\t)});
      \addplot[simBlue!70, samples=80, domain=0:360, variable=\t]
        ({sqrt(1.1)*cos(\t)}, {sqrt(0.4)*sin(\t)});
      \addplot[simBlue!45, samples=80, domain=0:360, variable=\t]
        ({sqrt(0.35)*cos(\t)}, {sqrt(0.13)*sin(\t)});
      \node[circle, fill=simRust, inner sep=1.3pt, label={[font=\scriptsize,text=simRust]below:$\vx^\star$}] at (axis cs:0,0) {};
      \coordinate (p) at (axis cs:1.35,0.78);
      \node[circle, fill=simInk, inner sep=1.1pt] at (p) {};
      \node[font=\scriptsize] at (axis cs:1.62,1.02) {$\vx$};
      \draw[flow, simRust] (p) -- (axis cs:0.55,0.18)
        node[midway, above right, font=\scriptsize, text=simRust] {$\vr=-\Grad\phi$};
    \end{axis}
  \end{scope}
\end{tikzpicture}
\caption[The quadratic bowl and its contours]{The function
$\phi(\vx)=\tfrac12\vx^{T}\mat{A}\vx-\vb^{T}\vx$ for an SPD $\mat{A}$. \emph{Left:}
the graph is an upward paraboloid---a bowl---with a single lowest point.
\emph{Right:} its level sets are concentric ellipses centered on the minimizer
$\vx^\star$, the solution of $\mat{A}\vx=\vb$. The ellipses are stretched along
$\mat{A}$'s eigenvectors by the reciprocals of its eigenvalues; the more
$\mat{A}$ stretches, the more eccentric they are. The residual $\vr=\vb-\mat{A}\vx
= -\Grad\phi$ points in the direction of steepest descent---straight down the
local slope, \emph{not}, in general, toward $\vx^\star$.}
\label{fig:bowl}
\end{figure}

\subsection{Why steepest descent zig-zags}
\label{sec:zigzag}

The most obvious way to descend a bowl is to repeatedly step in the steepest
downhill direction. That is \emph{steepest descent}: from the current $\vx$,
compute the residual $\vr=\vb-\mat{A}\vx$ (the downhill direction), and move
along it by the amount that minimizes $\phi$ on that line. The line-minimizing
step length has a clean closed form. Along $\vx+\alpha\vr$,
\[
  \phi(\vx+\alpha\vr) = \phi(\vx) - \alpha\,\vr^{T}\vr
    + \tfrac12\alpha^2\,\vr^{T}\mat{A}\vr,
\]
a parabola in $\alpha$ minimized at
\begin{equation}
\label{eq:sdstep}
  \alpha = \frac{\vr^{T}\vr}{\vr^{T}\mat{A}\vr}
         = \frac{\inner{\vr}{\vr}}{\inner{\mat{A}\vr}{\vr}} .
\end{equation}
Take the step, recompute the residual, repeat. It converges---$\phi$ decreases
every step and the bowl has a floor---but on a stretched bowl it converges
\emph{slowly}, and the reason is geometric and worth seeing.

When you minimize $\phi$ exactly along a line, you stop at the point where the
line is tangent to a level ellipse; there the new residual is
\emph{perpendicular to the line you just walked}. So each steepest-descent step
turns ninety degrees from the last. On a round bowl (eigenvalues all equal) the
first residual already points at the center and one step finishes. On a stretched
bowl (eigenvalues far apart) the downhill direction points mostly \emph{across}
the long valley rather than \emph{along} it, so the iterate crabs back and forth,
each step undoing part of the last, creeping toward the minimum in a long
zig-zag. \Cref{fig:zigzag} (left) is the classic picture. The narrower the
valley---the larger the ratio of largest to smallest eigenvalue---the worse the
crabbing. That ratio is the \emph{condition number}, and it will govern
everything.

\begin{keyidea}
Steepest descent stops, each step, where the new residual is perpendicular to the
direction just travelled. On a stretched (ill-conditioned) bowl, consecutive
perpendicular steps point mostly across the valley, so the iterate zig-zags and
crawls. The cure is not a better step \emph{length}---\cref{eq:sdstep} is already
optimal along its line---but a better set of \emph{directions}: directions chosen
so that progress along one is never spoiled by the next. Building those
directions is conjugate gradients.
\end{keyidea}

\section{Conjugacy: directions that do not interfere}
\label{sec:conjugacy}

The flaw in steepest descent is that minimizing along the new residual
\emph{re-introduces} error along the old direction: the steps are coupled. We
want a set of directions $\vp_0,\vp_1,\dots,\vp_{N-1}$ along which we can minimize
\emph{one at a time}, with the guarantee that optimizing along $\vp_k$ never
disturbs the optimization already done along $\vp_0,\dots,\vp_{k-1}$. The
condition that makes the directions independent in exactly this sense is
\emph{conjugacy}.

\begin{definition}[$\mat{A}$-conjugate directions]
\label{def:conjugate}
Two nonzero vectors $\vp_i,\vp_j$ are \emph{conjugate with respect to $\mat{A}$}
(or \emph{$\mat{A}$-orthogonal}) if
\[
  \inner{\mat{A}\vp_i}{\vp_j} = \vp_i^{T}\mat{A}\vp_j = 0,\qquad i\neq j .
\]
A set $\{\vp_k\}$ is $\mat{A}$-conjugate if its members are pairwise conjugate.
\end{definition}

Conjugacy is ordinary orthogonality measured in the \emph{energy inner product}
$\inner{\vu}{\vv}_{\mat{A}}:=\vu^{T}\mat{A}\vv$, which is a genuine inner product
precisely because $\mat{A}\succ0$. Its induced norm,
$\norm{\vu}_{\mat{A}}:=\sqrt{\vu^{T}\mat{A}\vu}$, is the \emph{energy norm}, and
it is the natural ruler for this problem: the quantity conjugate gradients
actually drives down is the energy-norm error, as \cref{sec:converge} will show.
Geometrically, $\mat{A}$-conjugate directions look skew in ordinary axes but are
perpendicular once the bowl is un-stretched into a round one. \Cref{fig:conjugate}
shows the two pictures side by side.

\begin{figure}[t]
\centering
\begin{tikzpicture}[font=\footnotesize]
  % --- left: skew in standard coords ---
  \begin{scope}
    \node[font=\small\bfseries, text=simBlue] at (0,2.55) {standard axes: skew};
    \begin{scope}[scale=0.95]
      \draw[simBlue, thick] (0,0) ellipse (1.9 and 0.95);
      \draw[simBlue!55] (0,0) ellipse (1.25 and 0.62);
      \node[circle, fill=simRust, inner sep=1.2pt] at (0,0) {};
      % two A-conjugate directions: skew here
      \draw[flow, simGreen, thick] (0,0) -- (1.7,0)
        node[right, font=\scriptsize, text=simGreen] {$\vp_0$};
      \draw[flow, simRust, thick] (0,0) -- (-0.62,0.78)
        node[above left, font=\scriptsize, text=simRust] {$\vp_1$};
      \node[font=\scriptsize, text=simGray] at (0.0,-1.45) {$\vp_0^{T}\mat{A}\vp_1=0$};
    \end{scope}
  \end{scope}
  \draw[simGray!50, thick] (3.7,-1.7) -- (3.7,2.7);
  % --- right: orthogonal in A-coords ---
  \begin{scope}[xshift=7.6cm]
    \node[font=\small\bfseries, text=simGreen] at (0,2.55) {energy ($\mat{A}$) axes: orthogonal};
    \begin{scope}[scale=0.95]
      \draw[simGreen, thick] (0,0) circle (1.25);
      \draw[simGreen!55] (0,0) circle (0.78);
      \node[circle, fill=simRust, inner sep=1.2pt] at (0,0) {};
      \draw[flow, simGreen, thick] (0,0) -- (1.25,0)
        node[right, font=\scriptsize, text=simGreen] {$\vp_0$};
      \draw[flow, simRust, thick] (0,0) -- (0,1.25)
        node[above, font=\scriptsize, text=simRust] {$\vp_1$};
      \draw[simGray] (0.22,0) -- (0.22,0.22) -- (0,0.22);
      \node[font=\scriptsize, text=simGray] at (0.0,-1.55) {perpendicular here};
    \end{scope}
  \end{scope}
\end{tikzpicture}
\caption[Conjugate directions in two coordinate systems]{$\mat{A}$-conjugate
directions, drawn twice. \emph{Left:} in ordinary coordinates the level sets are
ellipses and the two directions $\vp_0,\vp_1$ look \emph{skew}---they are not
perpendicular in the usual sense. \emph{Right:} change to coordinates that
un-stretch the bowl into a circle (the energy inner product
$\inner{\vu}{\vv}_{\mat{A}}=\vu^{T}\mat{A}\vv$), and the same two directions are
exactly \emph{perpendicular}. Conjugacy is orthogonality in the energy geometry.
This is why minimizing along $\vp_1$ cannot spoil the work done along $\vp_0$:
they are independent axes of the round bowl.}
\label{fig:conjugate}
\end{figure}

\subsection{Why conjugate directions finish in \texorpdfstring{$N$}{N} steps}
\label{sec:finite}

The payoff of conjugacy is a theorem that has no analogue for steepest descent:
with conjugate directions, exact minimization terminates after at most $N$ steps.
The idea is exactly the independence we asked for.

\begin{theorem}[Conjugate-direction termination]
\label{thm:finite}
Let $\{\vp_0,\dots,\vp_{N-1}\}$ be $\mat{A}$-conjugate (hence a basis of
$\Real^N$). Starting from any $\vx_0$, set
$\vx_{k+1}=\vx_k+\alpha_k\vp_k$ with $\alpha_k$ chosen to minimize $\phi$ along
$\vp_k$. Then $\vx_N=\vx^\star$, the exact solution of $\mat{A}\vx=\vb$.
\end{theorem}

\begin{proof}[Proof sketch]
Because the $\vp_k$ are a basis, write the initial error in them:
$\vx^\star-\vx_0=\sum_{k}\delta_k\vp_k$. Apply $\mat{A}$ and take the energy inner
product with $\vp_j$; conjugacy kills every cross term, leaving
$\delta_j=\inner{\vp_j}{\mat{A}(\vx^\star-\vx_0)}/\inner{\mat{A}\vp_j}{\vp_j}
=\inner{\vp_j}{\vr_0}/\inner{\mat{A}\vp_j}{\vp_j}$, where $\vr_0=\vb-\mat{A}\vx_0$.
Separately, the line-minimizing step length along $\vp_j$ works out (next
subsection) to exactly this same $\delta_j$. So the $j$-th step installs exactly
the $j$-th coordinate of the error and never touches the others---each conjugate
direction is corrected once, independently, and after $N$ corrections every
coordinate is right.
\end{proof}

This is the finite-termination property: \emph{in exact arithmetic, conjugate
gradients solves an $N\times N$ SPD system in at most $N$ steps}, because it lays
the $N$ independent corrections down one at a time with no interference. In
floating point the guarantee is softened (rounding spoils exact conjugacy), and
in practice we stop far earlier---after a residual tolerance is met, typically in
\emph{far fewer} than $N$ steps for a well-conditioned system, as
\cref{sec:converge} quantifies. But the finite-termination theorem is the
structural fact that distinguishes conjugate directions from steepest descent's
endless crawl.

\begin{keyidea}
Conjugate directions are \emph{$\mat{A}$-orthogonal}: $\vp_i^{T}\mat{A}\vp_j=0$
for $i\neq j$. Minimizing $\phi$ along each in turn corrects one independent
coordinate of the error and disturbs none of the others, so an $N$-dimensional
SPD system is solved exactly in at most $N$ such steps. This finite termination
is what steepest descent, with its coupled perpendicular steps, can never
promise.
\end{keyidea}

\section{Deriving the method}
\label{sec:derive}

We now have the goal---march along $\mat{A}$-conjugate directions---but not yet
the directions. A naive plan would precompute a full conjugate basis (a
Gram--Schmidt orthogonalization in the energy inner product, keeping every
previous direction). That costs $O(N^2)$ storage and is hopeless for the
million-unknown systems Palace solves. The miracle of conjugate gradients,
derived in this section, is that for an SPD operator the conjugate directions can
be built by a \emph{short recurrence}: each new direction needs only the current
residual and the \emph{immediately previous} direction. Nothing older is kept.

We build the method from three demands and read off each coefficient as the
unique value that satisfies them.

\subsection{The step length \texorpdfstring{$\alpha$}{alpha}}
\label{sec:alpha}

Suppose we have a current iterate $\vx_k$, its residual $\vr_k=\vb-\mat{A}\vx_k$,
and a search direction $\vp_k$. We move to $\vx_{k+1}=\vx_k+\alpha_k\vp_k$ and ask
for the $\alpha_k$ that minimizes $\phi$ along $\vp_k$---the line search of
\cref{sec:zigzag}, now along $\vp_k$ instead of along the residual. The same
parabola calculation gives
\begin{equation}
\label{eq:alpha}
  \alpha_k = \frac{\inner{\vr_k}{\vp_k}}{\inner{\mat{A}\vp_k}{\vp_k}} .
\end{equation}
Equivalently---and this is the form the algorithm uses---we will show below that
$\inner{\vr_k}{\vp_k}=\inner{\vr_k}{\vr_k}$, so
\begin{equation}
\label{eq:alpha2}
  \boxed{\;\alpha_k = \frac{\inner{\vr_k}{\vr_k}}{\inner{\mat{A}\vp_k}{\vp_k}}\;}
\end{equation}
The denominator $\inner{\mat{A}\vp_k}{\vp_k}=\norm{\vp_k}_{\mat{A}}^2$ is the
energy norm of the search direction, strictly positive because $\mat{A}\succ0$ and
$\vp_k\neq0$, so the division is always safe until the residual vanishes.
The numerator is just the squared length of the current residual---a quantity we
need anyway to test convergence. One operator apply $\mat{A}\vp_k$ and two inner
products give the step.

\subsection{The residual update}
\label{sec:residual}

Having stepped, we need the new residual $\vr_{k+1}=\vb-\mat{A}\vx_{k+1}$. Do not
recompute it from scratch with a fresh matrix-vector product---that would cost a
second apply. Instead subtract the change:
\[
  \vr_{k+1} = \vb-\mat{A}(\vx_k+\alpha_k\vp_k)
            = \vr_k - \alpha_k\,\mat{A}\vp_k .
\]
The product $\mat{A}\vp_k$ was already computed for the step length
\eqref{eq:alpha2}, so the residual update is a single \code{axpy} (a
scaled-vector add) and costs \emph{no} extra apply. This is the reason a conjugate
gradient iteration needs exactly \emph{one} operator application per step---the
governing cost metric of \cref{ch:krylov}. The recurrence
\begin{equation}
\label{eq:rupdate}
  \vr_{k+1} = \vr_k - \alpha_k\,\mat{A}\vp_k
\end{equation}
also drives two orthogonality facts we now harvest.

\subsection{The direction update and the short-recurrence miracle}
\label{sec:beta}

The new direction must be conjugate to all previous ones. We build it as the new
residual plus a correction along the \emph{previous} direction only:
\begin{equation}
\label{eq:pupdate}
  \vp_{k+1} = \vr_{k+1} + \beta_{k}\,\vp_k ,
\end{equation}
and choose $\beta_k$ to force conjugacy with $\vp_k$, i.e.
$\inner{\mat{A}\vp_{k+1}}{\vp_k}=0$. Substituting \eqref{eq:pupdate} and solving,
\[
  \beta_k = -\,\frac{\inner{\mat{A}\vr_{k+1}}{\vp_k}}{\inner{\mat{A}\vp_k}{\vp_k}} .
\]
This forces conjugacy against $\vp_k$. The astonishing fact---the heart of the
method---is that for an SPD operator, forcing conjugacy against the \emph{one}
previous direction automatically delivers conjugacy against \emph{all} the older
ones, for free. Two orthogonality relations make it happen, both consequences of
the residual recurrence \eqref{eq:rupdate}:

\begin{lemma}[CG orthogonalities]
\label{lem:orth}
Running the recurrences \eqref{eq:alpha2}, \eqref{eq:rupdate}, \eqref{eq:pupdate}
from $\vp_0=\vr_0$ produces, for all valid indices,
\begin{enumerate}[nosep, label=\textup{(\roman*)}]
  \item mutually orthogonal residuals: $\inner{\vr_i}{\vr_j}=0$ for $i\neq j$;
  \item $\mat{A}$-conjugate directions: $\inner{\mat{A}\vp_i}{\vp_j}=0$ for $i\neq j$.
\end{enumerate}
\end{lemma}

The two facts feed each other by induction: orthogonal residuals make the
direction built by \eqref{eq:pupdate} conjugate to \emph{every} prior direction
(not just the last), and conjugate directions in turn make the next residual
orthogonal to every prior residual. The full induction is in any of the further
reading; the point for us is its consequence. Because the residuals are mutually
orthogonal, the coefficient $\beta_k$ collapses to a ratio of \emph{consecutive
residual norms}. Using \eqref{eq:rupdate} to rewrite
$\mat{A}\vp_k=(\vr_k-\vr_{k+1})/\alpha_k$ and applying orthogonality term by term,
the messy $\beta_k$ above simplifies to
\begin{equation}
\label{eq:beta}
  \boxed{\;\beta_k = \frac{\inner{\vr_{k+1}}{\vr_{k+1}}}{\inner{\vr_k}{\vr_k}}\;}
\end{equation}
the celebrated Fletcher--Reeves form. (The same orthogonalities collapse
$\inner{\vr_k}{\vp_k}$ to $\inner{\vr_k}{\vr_k}$, which is the identity promised
in \cref{sec:alpha}.) No old direction, no old residual: $\beta_k$ is built from
the squared residual norm we just computed and the one from the step before.

\begin{keyidea}
\textbf{The short-recurrence miracle.} For an SPD operator, a direction made
conjugate to only the \emph{previous} direction is automatically conjugate to
\emph{all} previous directions. Consequently conjugate gradients keeps just two
vectors---the residual $\vr$ and the direction $\vp$---and two scalars---the step
length $\alpha$ and the coefficient $\beta=\norm{\vr_{k+1}}^2/\norm{\vr_k}^2$. It
never stores or re-orthogonalizes against the old directions, yet it builds a
full conjugate basis. This $O(N)$-memory short recurrence is what makes CG
practical at a million unknowns---and it is exactly what GMRES, on a
\emph{non}-symmetric operator, cannot do (\cref{ch:gmres}).
\end{keyidea}

\section{The algorithm}
\label{sec:algorithm}

Collecting \eqref{eq:alpha2}, \eqref{eq:rupdate}, \eqref{eq:pupdate},
\eqref{eq:beta} gives the whole method. We present it twice: first as ordinary
pseudocode, then as the one \code{iterate\_while} of the shared step kernel, to
make the tie to \cref{ch:fold} and \cref{ch:krylov} explicit.

\subsection{As pseudocode}

\begin{lstlisting}[style=l4style]
-- Conjugate gradients: solve A x = b for SPD A.
cg A b x0 tol =
  let r0 = b - A x0                    -- initial residual
      p0 = r0                          -- first direction is steepest descent
  in loop with x = x0, r = r0, p = p0:
       Ap    = A p                     -- THE one operator apply per step
       alpha = <r,r> / <Ap,p>          -- step length          (eq 26.6)
       x'    = x + alpha * p           -- advance the iterate
       r'    = r - alpha * Ap          -- update residual, no new apply (eq 26.7)
       if sqrt <r',r'> <= tol: return x'
       beta  = <r',r'> / <r,r>         -- direction coefficient (eq 26.11)
       p'    = r' + beta * p           -- new conjugate direction (eq 26.8)
       continue with x = x', r = r', p = p'
\end{lstlisting}

Per step: one apply $\mat{A}\vp$, two inner products ($\inner{\vr}{\vr}$ and
$\inner{\mat{A}\vp}{\vp}$; the next $\inner{\vr'}{\vr'}$ is reused as the
following step's numerator), and three scaled-vector adds (the updates of $\vx$,
$\vr$, and $\vp$). \Cref{fig:dataflow} draws this as a dataflow graph: one apply,
two dots, three \code{axpy}s. That fixed, tiny per-step recipe---independent of
$N$ beyond the vector lengths---is the whole computational kernel.

\begin{figure}[t]
\centering
\begin{tikzpicture}[font=\footnotesize, node distance=6mm and 9mm]
  \node[data, minimum width=9mm] (r) {$\vr$};
  \node[data, minimum width=9mm, below=10mm of r] (p) {$\vp$};
  \node[opbox, right=of p] (ap) {\code{A·p}\\\scriptsize 1 apply};
  \node[opbox, right=12mm of ap] (d2) {\code{<Ap,p>}\\\scriptsize dot};
  \node[opbox, above=8mm of d2] (d1) {\code{<r,r>}\\\scriptsize dot};
  \node[opbox, right=10mm of d2] (al) {$\alpha$};
  \node[opbox, right=12mm of d1] (xu) {\code{x += a p}\\\scriptsize axpy};
  \node[opbox, below=6mm of xu] (ru) {\code{r -= a Ap}\\\scriptsize axpy};
  \node[opbox, below=6mm of ru] (be) {$\beta=\frac{<r',r'>}{<r,r>}$};
  \node[opbox, right=10mm of ru] (pu) {\code{p = r'+b p}\\\scriptsize axpy};
  \node[product, right=9mm of pu] (out) {$\vx',\vr',\vp'$};
  \draw[flow] (r)-|(d1);
  \draw[flow] (p)--(ap);
  \draw[flow] (ap)--(d2);
  \draw[flow] (ap.east) .. controls +(0.5,-0.7) and +(-0.5,0) .. (ru.west);
  \draw[flow] (d1)--(al);  \draw[flow] (d2)--(al);
  \draw[flow] (al) .. controls +(0,0.9) and +(0,0.9) .. (xu.north);
  \draw[flow] (al) .. controls +(0.2,-0.3) and +(-0.2,0.3) .. (ru.east);
  \draw[flow] (p) .. controls +(-0.5,0) and +(-0.6,-0.4) .. (xu.west);
  \draw[flow] (xu)--(out);
  \draw[flow] (ru)--(be);  \draw[flow] (ru) .. controls +(0.3,0) and +(-0.3,0.2) .. (pu.west);
  \draw[flow] (be) .. controls +(0.6,0) and +(-0.3,-0.4) .. (pu.south);
  \draw[flow] (pu)--(out);
\end{tikzpicture}
\caption[The CG step dataflow]{One conjugate-gradient step as a dataflow graph:
\emph{one} operator apply $\mat{A}\vp$, \emph{two} inner products
($\inner{\vr}{\vr}$ and $\inner{\mat{A}\vp}{\vp}$), and \emph{three}
scaled-vector adds (\code{axpy}s) updating $\vx$, $\vr$, and $\vp$. The product
$\mat{A}\vp$ feeds both the denominator of $\alpha$ and the residual update---it
is computed once and reused, which is why a step costs a single apply. Compare the
threaded-carry picture of \cref{fig:carry}: this graph is the \emph{body} of one
\texttt{iterate\_while} step.}
\label{fig:dataflow}
\end{figure}

\subsection{As one \texttt{iterate\_while} of the step kernel}
\label{sec:iterate}

\Cref{ch:fold} taught that every iterative algorithm in this book is one
\code{iterate\_while}, and \cref{ch:krylov} taught that every Krylov solver is
that combinator folding a single \code{krylov\_step}. Conjugate gradients is the
canonical instance. We read the algorithm through the strata of \cref{ch:strata}.

The \emph{carry} is the run-time \code{SimState}---the iterate $\vx$, the
iteration count, the convergence flag, the residual statistics---exactly the
five-field record that is uniform across every Krylov method. The
\emph{ephemeral} workspace is the CG \code{Krylov} bundle: the residual $\vr$, the
direction $\vp$, the optional preconditioned residual $\vz$, and the two scalars
$\alpha,\beta$. It is solve-local scratch, born when the solve starts and
discarded when it returns, threaded as a plain value beside the carry rather than
through the monad. The \emph{predicate} reads the residual proxy folded into the
carry and stops at tolerance---the discipline of \cref{ch:fold}, never reading a
quantity the step has not yet produced. And the \emph{step} is one
\code{krylov\_step}: apply once, update the scalars, update the three vectors,
emit the residual norm as a demand-prunable output.

\begin{lstlisting}[style=l4style]
cg_solve :: OpParams -> SimState -> { final_state: SimState, trajectory: [...] }
cg_solve op s0 =
  iterate_while
    s0                                 -- carry: the SimState (x, it, converged, ...)
    (\s -> not s.converged && s.it < op.max_it)   -- predicate: pure on the carry
    (\s -> krylov_step op krylov s)     -- step: one apply, fold the CG body
-- where krylov_step performs, per call:
--   Ap    = apply_linop op.T  krylov.p          -- the ONE apply
--   alpha = dot(krylov.r, krylov.r) / dot(Ap, krylov.p)
--   x'    = axpy  alpha krylov.p    s.x          -- iterate update
--   r'    = axpy (-alpha) Ap        krylov.r     -- residual update
--   beta  = dot(r', r') / dot(krylov.r, krylov.r)
--   p'    = axpy  beta   krylov.p   r'           -- direction update
--   outputs = { residual_norm: sqrt (dot(r',r')) }   -- demand-prunable
\end{lstlisting}

The coefficient $\beta=\norm{\vr'}^2/\norm{\vr}^2$ reads the previous step's
residual-norm-squared. That single backward dependence is the
\emph{recurrence carry} of \cref{ch:strata}'s fourth stratum---the scalar
threaded from one inner step to the next and reborn at each solve. In the
branch-in-body form (\cref{ch:fold}'s Form~A) it lives as a field of the workspace
with an \lstinline[style=l4style]|if it == 0| guard for the first step (where
$\vp_0=\vr_0$, no previous direction); the first-iteration-unrolling rotation
splits it into a bootstrap step and a branch-free steady step
(\code{iterate\_while\_with\_prev}), exactly as \cref{ch:fold} described.
\Cref{fig:carry} draws the CG carry threading forward with the residual history
peeling off into the trajectory.

\begin{figure}[t]
\centering
\begin{tikzpicture}[font=\footnotesize, node distance=6mm]
  \node[data, minimum width=14mm, align=center] (s0) at (0,0) {\code{SimState}\\$\vx_0$};
  \node[diamond, draw=simBlue, fill=simBgBlue, aspect=1.7, inner sep=1pt,
        font=\scriptsize] (p0) at (2.2,0) {res$>$tol?};
  \node[opbox, minimum width=15mm, align=center] (k0) at (4.6,0) {\code{krylov\_step}\\\scriptsize 1 apply};
  \node[data, minimum width=14mm, align=center] (s1) at (7.2,0) {\code{SimState}\\$\vx_1$};
  \node[diamond, draw=simBlue, fill=simBgBlue, aspect=1.7, inner sep=1pt,
        font=\scriptsize] (p1) at (9.4,0) {res$>$tol?};
  \node[font=\scriptsize, simGray] (more) at (11.0,0) {$\cdots$};
  \draw[flow] (s0)--(p0);
  \draw[flow] (p0)-- node[above,font=\scriptsize]{yes} (k0);
  \draw[flow] (k0)-- node[above,font=\scriptsize]{\code{.state}} (s1);
  \draw[flow] (s1)--(p1);
  \draw[flow] (p1)-- node[above,font=\scriptsize]{yes} (more);
  % ephemeral Krylov threaded below as plain value
  \node[data, fill=simBgGreen, draw=simGreen, minimum width=20mm] (kr0) at (4.6,-1.5)
    {\scriptsize\code{Krylov}: $\vr,\vp,\alpha,\beta$};
  \draw[refedge, simGreen] (kr0.north) -- (k0.south);
  \draw[backflow, simGreen] (k0.south east) .. controls +(0.6,-0.8) and +(0.6,0.4) .. (kr0.east)
    node[midway, right=2mm, font=\scriptsize, text=simGreen] {plain value};
  % trajectory of residual norms peeling off
  \node[data, fill=simBgGold, draw=simGold, minimum width=40mm] (tj) at (6.5,-2.7)
    {\code{trajectory}: $[\,\norm{\vr_0},\ \norm{\vr_1},\ \dots\,]$};
  \draw[refedge, simGold!80!black] (k0.south) |- (tj.west);
  % converged exit
  \node[product] (fin) at (9.4,-1.5) {\code{final\_state}: $\vx^\star$};
  \draw[backflow] (p1.south)--(fin.north);
\end{tikzpicture}
\caption[CG as one \texttt{iterate\_while}]{Conjugate gradients as a single
\texttt{iterate\_while} (\cref{ch:fold}) folding one \code{krylov\_step}
(\cref{ch:krylov}). The carry is the run-time \code{SimState}
(\cref{ch:strata})---iterate $\vx$, counter, flags---threaded along the top; the
predicate reads the residual proxy folded into the carry and exits at tolerance.
The ephemeral \code{Krylov} workspace ($\vr,\vp,\alpha,\beta$) is threaded
\emph{beside} the carry as a plain value (green), reborn at solve entry and
discarded at return. Each step's residual norm peels off into the
\code{trajectory} (gold)---computed only if a consumer reads it, by the
demand-pruning of \cref{ch:fold}.}
\label{fig:carry}
\end{figure}

This is the structural payoff of Parts~II and~III landing on a real solver:
nothing in \cref{sec:algorithm} is new control structure. CG is the fold of
\cref{ch:fold}, over the strata of \cref{ch:strata}, applying the operator of
\cref{ch:krylov} once per step. What is special to CG lives entirely in the
\emph{body}---the four coefficient formulas we derived---not in the loop.

\section{Convergence}
\label{sec:converge}

Finite termination in $N$ steps (\cref{thm:finite}) is reassuring but useless as a
performance promise: $N$ is in the millions, and rounding spoils exact conjugacy
long before then. What makes CG \emph{fast} is that it usually drives the error
below tolerance in \emph{far fewer} steps---and how many fewer is governed by one
number, the condition number.

\begin{definition}[Condition number]
\label{def:kappa}
For an SPD operator $\mat{A}$ with eigenvalues
$0<\lambda_{\min}=\lambda_1\le\cdots\le\lambda_N=\lambda_{\max}$, the (spectral)
\emph{condition number} is
$\kappa=\kappa(\mat{A})=\lambda_{\max}/\lambda_{\min}$.
\end{definition}

$\kappa$ is the eccentricity of the bowl: $\kappa=1$ is a perfect round bowl (one
step finishes), and large $\kappa$ is a long narrow valley (the steepest-descent
zig-zag of \cref{sec:zigzag}). The central convergence result bounds the
energy-norm error of the $k$-th CG iterate by a power that shrinks as $\kappa$
shrinks.

\begin{theorem}[CG energy-norm convergence]
\label{thm:converge}
Let $\vx^\star$ solve $\mat{A}\vx=\vb$ with $\mat{A}=\mat{A}^{T}\succ0$ and
$\kappa=\kappa(\mat{A})$. The $k$-th conjugate-gradient iterate satisfies
\begin{equation}
\label{eq:bound}
  \norm{\vx^\star-\vx_k}_{\mat{A}}
  \;\le\; 2\!\left(\frac{\sqrt{\kappa}-1}{\sqrt{\kappa}+1}\right)^{\!k}
  \norm{\vx^\star-\vx_0}_{\mat{A}} .
\end{equation}
\end{theorem}

The structure of the proof is worth a sentence even though we do not give it in
full: after $k$ steps the CG iterate is the \emph{best possible} approximation, in
the energy norm, drawn from the degree-$k$ Krylov subspace
$\operatorname{span}\{\vr_0,\mat{A}\vr_0,\dots,\mat{A}^{k}\vr_0\}$
(\cref{ch:krylov})---that optimality is what conjugacy buys. Bounding the error
then becomes a question about polynomials: how small can a degree-$k$ polynomial
$q$ with $q(0)=1$ be made across the eigenvalue interval
$[\lambda_{\min},\lambda_{\max}]$? The answer is given by Chebyshev polynomials,
and it produces exactly the factor in \eqref{eq:bound}. The
$\sqrt{\kappa}$---not $\kappa$---is the decisive feature: CG converges at the rate
set by the \emph{square root} of the condition number, which is why it so
dramatically outperforms steepest descent (whose analogous bound carries $\kappa$
itself).

\begin{figure}[t]
\centering
\begin{tikzpicture}
  \begin{semilogyaxis}[
      width=12.2cm, height=7.0cm,
      xlabel={iteration $k$}, ylabel={relative energy-norm error (bound)},
      xmin=0, xmax=60, ymin=1e-8, ymax=3,
      grid=both, grid style={simGray!18},
      legend pos=north east, legend cell align=left,
      legend style={font=\footnotesize, draw=simGray!40},
      label style={font=\footnotesize}, tick label style={font=\scriptsize},
      every axis plot/.append style={very thick}]
    % bound 2*c^k with c=(sqrt(kappa)-1)/(sqrt(kappa)+1); written as 2*exp(k*ln c)
    \addplot[simGreen, domain=0:60, samples=61]
      {2*exp(x*ln((sqrt(10)-1)/(sqrt(10)+1)))};
    \addlegendentry{$\kappa=10$ (well-conditioned)}
    \addplot[simBlue, domain=0:60, samples=61]
      {2*exp(x*ln((sqrt(100)-1)/(sqrt(100)+1)))};
    \addlegendentry{$\kappa=100$}
    \addplot[simGold, domain=0:60, samples=61]
      {2*exp(x*ln((sqrt(1000)-1)/(sqrt(1000)+1)))};
    \addlegendentry{$\kappa=1000$}
    \addplot[simRust, domain=0:60, samples=61]
      {2*exp(x*ln((sqrt(10000)-1)/(sqrt(10000)+1)))};
    \addlegendentry{$\kappa=10000$ (ill-conditioned)}
    \draw[simGray, densely dashed] (axis cs:0,1e-6)--(axis cs:60,1e-6)
      node[pos=0.07, above, font=\scriptsize, text=simGray] {tol $=10^{-6}$};
  \end{semilogyaxis}
\end{tikzpicture}
\caption[CG convergence versus condition number]{The CG energy-norm error bound
$2\big((\sqrt\kappa-1)/(\sqrt\kappa+1)\big)^k$ of \cref{thm:converge}, for four
condition numbers, on a log scale. A well-conditioned system ($\kappa=10$) reaches
a $10^{-6}$ tolerance in a handful of steps; each tenfold increase in $\kappa$
roughly triples the iteration count, because the rate depends on $\sqrt\kappa$.
This is the entire motivation for preconditioning (\cref{ch:precond}): replace
$\mat{A}$ by $\mat{M}^{-1}\mat{A}$ with a far smaller condition number and slide
down to the green curve. The dependence on $\sqrt\kappa$ rather than $\kappa$ is
CG's decisive advantage over steepest descent.}
\label{fig:convplot}
\end{figure}

\Cref{fig:convplot} plots the bound for several condition numbers. The lesson is
stark: a well-conditioned system converges in a handful of steps; a poorly
conditioned one crawls. Since the operators Palace assembles are routinely
ill-conditioned---fine meshes and material contrasts push $\kappa$ into the
thousands or worse---raw CG is too slow, and the fix is to \emph{change the
condition number}. That is preconditioning, the subject of \cref{ch:precond} and
the reason \cref{sec:pcg} exists.

\subsection{Clustered eigenvalues and superlinear convergence}
\label{sec:cluster}

The bound \eqref{eq:bound} is worst-case; CG often does much better, and the
reason is the polynomial picture again. The error after $k$ steps is small if
\emph{some} degree-$k$ polynomial with $q(0)=1$ is small on the spectrum---and a
polynomial can be made small on a few tight \emph{clusters} of eigenvalues using
far fewer degrees than it would need to be small across a whole wide interval.
Concretely, if the spectrum consists of $m$ tight clusters, CG behaves as if
$\kappa$ were the spread \emph{within} a cluster, and it ``uses up'' the
outliers a few per step. The visible effect is \emph{superlinear convergence}: the
residual drops slowly at first, then accelerates as CG eliminates the extreme
eigenvalues and is left working only against the cluster. This is also why a good
preconditioner---one that clusters the spectrum near $1$---helps so much more than
the bare $\kappa$ improvement suggests.

\section{Preconditioned conjugate gradients}
\label{sec:pcg}

If convergence is governed by $\kappa(\mat{A})$, the way to make CG fast is to
solve a system with a smaller condition number that has the same solution.
Preconditioning supplies a \emph{preconditioner} $\mat{M}\approx\mat{A}$, cheap to
invert, such that $\mat{M}^{-1}\mat{A}$ is far better conditioned than $\mat{A}$.
The full theory is \cref{ch:precond}; here we record only how the CG \emph{loop}
changes, because the change is small and structural.

The naive idea---run CG on $\mat{M}^{-1}\mat{A}\vx=\mat{M}^{-1}\vb$---breaks the
SPD requirement, since $\mat{M}^{-1}\mat{A}$ is generally not symmetric. The
preconditioned conjugate gradient (PCG) algorithm finesses this by running CG in
the inner product defined by $\mat{M}$, which keeps every operator effectively
symmetric and recovers an honest CG. The upshot is the same recurrence with one
extra vector and one extra apply per step: at each step, form the
\emph{preconditioned residual} $\vz=\mat{M}^{-1}\vr$, and replace the residual
inner products $\inner{\vr}{\vr}$ in $\alpha$ and $\beta$ by the
\emph{$\mat{M}$-weighted} inner products $\inner{\vz}{\vr}$:
\begin{align}
\label{eq:pcg-alpha}
  \vz_k &= \mat{M}^{-1}\vr_k, &
  \alpha_k &= \frac{\inner{\vz_k}{\vr_k}}{\inner{\mat{A}\vp_k}{\vp_k}}, \\
\label{eq:pcg-beta}
  \beta_k &= \frac{\inner{\vz_{k+1}}{\vr_{k+1}}}{\inner{\vz_k}{\vr_k}}, &
  \vp_{k+1} &= \vz_{k+1} + \beta_k\,\vp_k .
\end{align}
When $\mat{M}=\mat{I}$ this is exactly the unpreconditioned method of
\cref{sec:algorithm}---$\vz=\vr$ and the formulas collapse back. The cost is one
preconditioner apply $\vz=\mat{M}^{-1}\vr$ per step, on top of the one operator
apply, and the storage of the extra vector $\vz$ (the optional \code{z?} field of
the CG \code{Krylov} bundle in \cref{ch:strata}). \Cref{fig:pcg} overlays the
extra apply on the dataflow of \cref{fig:dataflow}.

\begin{figure}[t]
\centering
\begin{tikzpicture}[font=\footnotesize, node distance=7mm and 11mm]
  \node[data, minimum width=9mm] (r) {$\vr$};
  \node[extern, right=of r] (M) {$\vz=\mat{M}^{-1}\vr$\\\scriptsize +1 apply};
  \node[opbox, right=of M] (dot) {\code{<z,r>}\\\scriptsize dot};
  \node[data, minimum width=9mm, below=12mm of r] (p) {$\vp$};
  \node[opbox, right=of p] (ap) {\code{A·p}\\\scriptsize 1 apply};
  \node[opbox, right=of ap] (dap) {\code{<Ap,p>}};
  \node[opbox, right=10mm of dap] (al) {$\alpha,\beta$};
  \node[opbox, right=10mm of dot] (upd) {\code{x,r,p}\\updates\\\scriptsize 3 axpy};
  \node[product, below=10mm of upd] (out) {$\vx',\vr',\vp'$};
  \draw[flow] (r)--(M); \draw[flow] (M)--(dot);
  \draw[flow] (p)--(ap); \draw[flow] (ap)--(dap);
  \draw[flow] (dot)--(al); \draw[flow] (dap)--(al);
  \draw[flow] (al) .. controls +(0,0.6) and +(0.4,-0.2) .. (upd.east);
  \draw[flow] (dot) .. controls +(0.4,0.3) and +(-0.4,0.2) .. (upd.west);
  \draw[flow] (M.south) .. controls +(0,-0.5) and +(-0.6,0.4) .. (upd.south west)
    node[midway, below left=-1mm, font=\scriptsize, text=simViolet] {$\vz$ into $\vp'$};
  \draw[flow] (upd)--(out);
\end{tikzpicture}
\caption[Preconditioned CG: the extra apply]{Preconditioned conjugate gradients
adds exactly one operation to the loop of \cref{fig:dataflow}: a preconditioner
apply $\vz=\mat{M}^{-1}\vr$ (violet, dashed---an \emph{opaque} operator surface,
\cref{ch:krylov}). The residual inner products $\inner{\vr}{\vr}$ become the
$\mat{M}$-weighted products $\inner{\vz}{\vr}$ in both $\alpha$ and $\beta$, and
the direction is built from $\vz$ rather than $\vr$. Everything else---the single
operator apply $\mat{A}\vp$, the three \code{axpy}s, the short
recurrence---is unchanged. PCG is CG with one extra apply and one extra vector.}
\label{fig:pcg}
\end{figure}

\begin{pitfall}
\textbf{Do not run CG on a non-SPD system.} Every guarantee in this
chapter---finite termination, the energy-norm bound, even the line-minimizing step
length being a \emph{minimum} rather than a saddle---rests on
$\mat{A}=\mat{A}^{T}\succ0$ (\cref{def:spd}). Hand CG a non-symmetric operator and
the energy ``inner product'' $\vu^{T}\mat{A}\vv$ is not symmetric, conjugacy is
ill-defined, and the iteration can stall or diverge with no warning. Hand it a
symmetric but \emph{indefinite} operator (negative eigenvalues---the
time-harmonic Maxwell operator $\mat{K}-\omega^2\mat{M}$ above resonance is one)
and the denominator $\inner{\mat{A}\vp}{\vp}$ can vanish or go negative,
\emph{dividing by zero or stepping uphill}. The symptom is a residual that
refuses to fall or suddenly blows up. The fix is to use the right method: GMRES
(\cref{ch:gmres}) for non-symmetric systems, MINRES for symmetric-indefinite ones.
A subtle trap: applying a non-SPD \emph{preconditioner} $\mat{M}$ inside PCG
breaks SPD-ness just as surely as a non-SPD $\mat{A}$---the preconditioner must
itself be SPD for \eqref{eq:pcg-alpha}--\eqref{eq:pcg-beta} to define an honest CG.
\end{pitfall}

\section{The load-bearing numerical note: reduction order}
\label{sec:numerics}

Every coefficient in CG---$\alpha$, $\beta$, the convergence test---is built from
\emph{inner products}, $\inner{\vu}{\vv}=\sum_i u_i v_i$. An inner product is a
\emph{global reduction}: it sums a product over all $N$ entries of the vectors.
And floating-point addition is \emph{not associative}: $(a+b)+c$ and $a+(b+c)$ can
differ in their last bits. So the value of an inner product---and therefore the
exact $\alpha$, the exact $\beta$, the exact iterate $\vx_k$, and the exact
iteration at which the residual first dips below tolerance---depends on the
\emph{order} in which the reduction is summed.

This is not a defect to be engineered away; it is a load-bearing property in the
sense of \cref{ch:bigidea}. Different reduction orders (a different summation tree,
a different thread count, a different device) give bit-different but equally valid
CG runs: the mathematics is identical, the floating-point realization is not. Two
runs of the ``same'' solve on two machines may take a different number of
iterations and return iterates differing in their last bits. The converged answer
is correct to tolerance either way; what is \emph{not} reproducible bit-for-bit is
the trajectory. The Palace source records this explicitly, and we treat the
reduction-order semantics in full in \cref{app:numerics}. The discipline for now is
simply to \emph{name} it: an inner product is a reduction, a reduction has an
order, and the order is observable in the last bits. When a downstream consumer
demands bit-reproducibility, it must pin the reduction order---CG does not, and
cannot, promise it for free.

\section{Worked examples}

\begin{workedexample}{CG on a $2\times2$ SPD system: two steps to exact}{2x2}
We run unpreconditioned CG on a $2\times2$ SPD system and watch it reach the exact
answer in two steps, as \cref{thm:finite} promises for $N=2$. Take
\[
  \mat{A}=\begin{pmatrix}4 & 1\\[1pt] 1 & 3\end{pmatrix},\qquad
  \vb=\begin{pmatrix}1\\[1pt] 2\end{pmatrix},\qquad
  \vx_0=\begin{pmatrix}0\\[1pt] 0\end{pmatrix}.
\]
$\mat{A}$ is symmetric, and its eigenvalues
$\tfrac{7\pm\sqrt5}{2}\approx\{2.38,\,4.62\}$ are positive, so it is SPD. The
exact solution is $\vx^\star=\mat{A}^{-1}\vb=\tfrac1{11}(1,7)^{T}
\approx(0.0909,\,0.6364)^{T}$.

\tcblower
\textbf{Step 0.} Residual $\vr_0=\vb-\mat{A}\vx_0=(1,2)^{T}$; first direction
$\vp_0=\vr_0=(1,2)^{T}$.
\begin{align*}
  \mat{A}\vp_0 &= (4\cdot1+1\cdot2,\;1\cdot1+3\cdot2)^{T}=(6,7)^{T}, &
  \inner{\vr_0}{\vr_0} &= 1+4 = 5,\\
  \inner{\mat{A}\vp_0}{\vp_0} &= 6\cdot1+7\cdot2 = 20, &
  \alpha_0 &= \tfrac{5}{20} = 0.25.
\end{align*}
Update: $\vx_1=\vx_0+\alpha_0\vp_0=(0.25,\,0.5)^{T}$ and
$\vr_1=\vr_0-\alpha_0\mat{A}\vp_0=(1,2)^{T}-0.25(6,7)^{T}=(-0.5,\,0.25)^{T}$.
Check orthogonality (\cref{lem:orth}): $\inner{\vr_1}{\vr_0}=-0.5+0.5=0$.~\checkmark

\textbf{Direction update.}
$\inner{\vr_1}{\vr_1}=0.25+0.0625=0.3125$, so
$\beta_0=\tfrac{0.3125}{5}=0.0625$, and
$\vp_1=\vr_1+\beta_0\vp_0=(-0.5,0.25)^{T}+0.0625(1,2)^{T}=(-0.4375,\,0.375)^{T}$.
Check conjugacy: $\mat{A}\vp_1=(-1.375,\,0.6875)^{T}$, and
$\inner{\mat{A}\vp_1}{\vp_0}=-1.375+1.375=0$.~\checkmark

\textbf{Step 1.}
$\inner{\mat{A}\vp_1}{\vp_1}=(-1.375)(-0.4375)+(0.6875)(0.375)=0.8594$, so
$\alpha_1=\tfrac{0.3125}{0.8594}=0.3636$.
\[
  \vx_2 = \vx_1+\alpha_1\vp_1
        = (0.25,0.5)^{T}+0.3636\,(-0.4375,0.375)^{T}
        = (0.0909,\,0.6364)^{T}.
\]
This is $\vx^\star$ to four places: in exactly $N=2$ steps the residual is zero
and CG has landed on the solution. The two steps installed the two conjugate
coordinates of the error, one each, with no interference---finite termination made
arithmetic.
\end{workedexample}

\begin{workedexample}{Steepest descent versus CG on an ill-conditioned bowl}{zigzag}
We put steepest descent and CG on the same stretched system and watch one
zig-zag while the other finishes in two steps. Take the diagonal SPD operator
\[
  \mat{A}=\begin{pmatrix}10 & 0\\[1pt] 0 & 1\end{pmatrix},\qquad
  \vb=\begin{pmatrix}0\\[1pt] 0\end{pmatrix},\qquad
  \vx_0=\begin{pmatrix}1\\[1pt] 1\end{pmatrix},
\]
so $\vx^\star=0$ and $\phi(\vx)=5x_1^2+\tfrac12x_2^2$ is a valley ten times
steeper across than along: $\kappa=10$. \Cref{fig:zigzag} plots both iterate
paths on the contours.

\tcblower
\textbf{Steepest descent} (\cref{eq:sdstep}). From $\vx_0=(1,1)^{T}$,
$\vr_0=-\mat{A}\vx_0=(-10,-1)^{T}$,
$\alpha_0=\inner{\vr_0}{\vr_0}/\inner{\mat{A}\vr_0}{\vr_0}
=\tfrac{101}{1010}=0.1$, giving $\vx_1=(0,\,0.9)^{T}$. The next residual
$(0,-0.9)^{T}$ is \emph{perpendicular} to the last step, so the iterate turns
ninety degrees and lands at $(0.818,0)^{T}$, then turns again. It crabs
$(1,1)\to(0,0.9)\to(0.82,0)\to(0,0.74)\to\cdots$, the error shrinking by only a
factor $\approx0.82$ per step---dozens of steps to reach $10^{-6}$.

\textbf{Conjugate gradients}. The same first step gives $\vx_1=(0,0.9)^{T}$ and
$\vr_1=(0,-0.9)^{T}$. But now $\beta_0=\inner{\vr_1}{\vr_1}/\inner{\vr_0}{\vr_0}
=\tfrac{0.81}{101}=0.00802$, and the direction
$\vp_1=\vr_1+\beta_0\vp_0=(-0.0802,\,-0.908)^{T}$ \emph{bends away} from the pure
residual so as to be $\mat{A}$-conjugate to $\vp_0$. One step along it,
$\alpha_1=0.901$, lands exactly on $\vx^\star=0$. \textbf{Two steps, done}---the
same finite termination as \cref{wex:2x2}, now visibly beating the
steepest-descent crawl on the identical bowl.
\end{workedexample}

\begin{figure}[t]
\centering
\begin{tikzpicture}
  \begin{axis}[
      width=12.4cm, height=7.2cm, axis equal image,
      xlabel={$x_1$}, ylabel={$x_2$},
      xmin=-0.55, xmax=1.15, ymin=-0.35, ymax=1.2,
      label style={font=\footnotesize}, tick label style={font=\scriptsize},
      legend pos=north west, legend cell align=left,
      legend style={font=\footnotesize, draw=simGray!40}]
    % elliptical contours of phi = 5 x1^2 + 0.5 x2^2 (levels 0.5, 2, 5)
    \addplot[simGray!45, samples=80, domain=0:360, variable=\t, forget plot]
      ({sqrt(0.5/5)*cos(\t)}, {sqrt(0.5/0.5)*sin(\t)});
    \addplot[simGray!45, samples=80, domain=0:360, variable=\t, forget plot]
      ({sqrt(2/5)*cos(\t)}, {sqrt(2/0.5)*sin(\t)});
    \addplot[simGray!45, samples=80, domain=0:360, variable=\t, forget plot]
      ({sqrt(5/5)*cos(\t)}, {sqrt(5/0.5)*sin(\t)});
    % steepest descent zig-zag path
    \addplot[simRust, very thick, mark=*, mark size=1.3pt]
      coordinates {(1,1) (0,0.9) (0.818,0) (0,0.744) (0.676,0) (0,0.614)};
    \addlegendentry{steepest descent (zig-zag)}
    % CG path: 2 steps
    \addplot[simBlue, very thick, mark=square*, mark size=1.6pt]
      coordinates {(1,1) (0,0.9) (0,0)};
    \addlegendentry{conjugate gradients (2 steps)}
    \node[circle, fill=simGreen, inner sep=1.4pt,
      label={[font=\scriptsize,text=simGreen]below right:$\vx^\star$}] at (axis cs:0,0) {};
  \end{axis}
\end{tikzpicture}
\caption[Steepest descent versus CG on the contours]{The classic comparison, on
the bowl $\phi=5x_1^2+\tfrac12 x_2^2$ ($\kappa=10$) of \cref{wex:zigzag}.
\emph{Steepest descent} (rust) takes the locally steepest step each time and so
turns ninety degrees at every iterate, zig-zagging down the valley and needing
many steps. \emph{Conjugate gradients} (blue) takes the same first step, then
bends its second direction to be $\mat{A}$-conjugate to the first and lands
exactly on $\vx^\star$ (green) in \emph{two} steps---the finite termination of
\cref{thm:finite}. The conjugacy correction is precisely the difference between
the crawl and the finish.}
\label{fig:zigzag}
\end{figure}

\begin{workedexample}{How many iterations for a target tolerance?}{count}
The bound \eqref{eq:bound} lets us \emph{predict} the iteration count before
running, given the condition number. Suppose an electrostatic solve assembles an
operator with $\kappa=500$ and we want the relative energy-norm error below
$\varepsilon=10^{-6}$. How many CG iterations does \cref{thm:converge} guarantee?

\tcblower
We need $k$ with $2\,c^{k}\le\varepsilon$, where
$c=\dfrac{\sqrt\kappa-1}{\sqrt\kappa+1}$. With $\kappa=500$,
$\sqrt\kappa=22.36$, so
\[
  c = \frac{22.36-1}{22.36+1} = \frac{21.36}{23.36} = 0.9144 .
\]
Solving $2\,c^{k}\le10^{-6}$:
\[
  k \;\ge\; \frac{\ln(\varepsilon/2)}{\ln c}
    = \frac{\ln(5\times10^{-7})}{\ln(0.9144)}
    = \frac{-14.51}{-0.0895}
    \approx 162 .
\]
So the bound guarantees $10^{-6}$ within about \textbf{162 iterations}. Now
precondition down to $\kappa=10$ (a realistic multigrid target,
\cref{ch:precond}): $\sqrt\kappa=3.162$,
$c=\tfrac{2.162}{4.162}=0.5195$, and
\[
  k \ge \frac{-14.51}{\ln(0.5195)} = \frac{-14.51}{-0.6548}\approx 23 .
\]
\textbf{162 iterations become 23}---a sevenfold speedup, each preconditioned
iteration only modestly more expensive (one extra apply, \cref{sec:pcg}). And
because the bound is worst-case, clustered spectra (\cref{sec:cluster}) routinely
beat even these counts. This single calculation---iterations scale like
$\sqrt\kappa$---is the entire economic case for the preconditioning of
\cref{ch:precond}.
\end{workedexample}

\summarysection
Conjugate gradients solves $\mat{A}\vx=\vb$ for a \emph{symmetric positive
definite} operator $\mat{A}=\mat{A}^{T}\succ0$---the electrostatic, magnetostatic,
and inner-Poisson systems of this book. Solving is minimizing the quadratic
$\phi(\vx)=\tfrac12\vx^{T}\mat{A}\vx-\vb^{T}\vx$, whose level sets are ellipses
and whose negative gradient is the residual $\vr=\vb-\mat{A}\vx$. Steepest
descent takes the locally steepest step and zig-zags across an ill-conditioned
valley; CG instead marches along \emph{$\mat{A}$-conjugate} directions
($\vp_i^{T}\mat{A}\vp_j=0$), which correct independent error coordinates and reach
the exact answer in at most $N$ steps. The method is built from four formulas: a
step length $\alpha_k=\inner{\vr_k}{\vr_k}/\inner{\mat{A}\vp_k}{\vp_k}$, a residual
update $\vr_{k+1}=\vr_k-\alpha_k\mat{A}\vp_k$ (no second apply), a direction update
$\vp_{k+1}=\vr_{k+1}+\beta_k\vp_k$, and the coefficient
$\beta_k=\inner{\vr_{k+1}}{\vr_{k+1}}/\inner{\vr_k}{\vr_k}$. The
\emph{short-recurrence miracle}---conjugacy against the previous direction yields
conjugacy against all---keeps the cost at \emph{one apply, two dots, three axpys}
per step and the memory at $O(N)$. As a program, CG is one \code{iterate\_while}
(\cref{ch:fold}) folding one \code{krylov\_step} (\cref{ch:krylov}): the carry is
the run-time \code{SimState}, the ephemeral $\vr,\vp$ are the \code{Krylov}
workspace, and $\beta$'s backward dependence is the recurrence carry
(\cref{ch:strata}). Convergence is governed by the condition number $\kappa$: the
energy-norm error shrinks like
$\big((\sqrt\kappa-1)/(\sqrt\kappa+1)\big)^k$, so well-conditioned (or clustered)
spectra converge fast and ill-conditioned ones crawl---the motivation for
preconditioning (\cref{ch:precond}), which PCG folds in as one extra apply per
step. Finally, the inner products are global reductions whose summation order
changes the last bits, a load-bearing non-determinism (\cref{app:numerics}).

\furtherreading
The definitive textbook account of conjugate gradients, the Krylov-subspace
optimality behind the convergence bound, and preconditioned CG is
Saad~\citep{saad2003}; the polynomial (Chebyshev) argument for
\cref{thm:converge} and the clustered-spectrum superlinear story are developed
there in full. Trefethen and Bau~\citep{trefethenbau1997} give the cleanest short
derivation of CG as the Krylov-subspace energy-norm minimizer and the most
geometric treatment of the convergence rate's $\sqrt\kappa$ dependence; their
Lectures~38 and~40 pair naturally with \cref{sec:converge}. For the numerical
linear algebra---the orthogonality lemmas of \cref{lem:orth}, the stability of the
recurrences in floating point, and the relationship to the Lanczos process of
\cref{ch:krylovschur}---Golub and Van Loan~\citep{golubvanloan2013} is the
standard reference.

\exercisessection
\begin{enumerate}[nosep]
  \item Verify directly that $\Grad\phi(\vx)=\mat{A}\vx-\vb$ for
  $\phi(\vx)=\tfrac12\vx^{T}\mat{A}\vx-\vb^{T}\vx$ when $\mat{A}=\mat{A}^{T}$.
  Where exactly does symmetry get used? Show that without symmetry the gradient is
  $\tfrac12(\mat{A}+\mat{A}^{T})\vx-\vb$, and explain why CG therefore ``sees'' only
  the symmetric part of a non-symmetric operator.
  \item Run CG by hand on
  $\mat{A}=\bigl(\begin{smallmatrix}2&0\\0&5\end{smallmatrix}\bigr)$,
  $\vb=(2,5)^{T}$, $\vx_0=0$. How many steps does it take, and why fewer than the
  generic $N=2$? (Relate your answer to the eigenvalue structure and
  \cref{sec:cluster}.)
  \item Derive the steepest-descent step length \eqref{eq:sdstep} by minimizing
  $\phi(\vx+\alpha\vr)$ over $\alpha$. Then show the next residual satisfies
  $\inner{\vr_{k+1}}{\vr_k}=0$, the perpendicularity that causes the zig-zag.
  \item Show that the line-minimizing step length along $\vp_k$ is
  \eqref{eq:alpha}, and then, assuming the orthogonality
  $\inner{\vr_k}{\vp_{k-1}}=0$ and the direction update, prove
  $\inner{\vr_k}{\vp_k}=\inner{\vr_k}{\vr_k}$, so \eqref{eq:alpha} becomes the
  algorithm's form \eqref{eq:alpha2}.
  \item Using the bound \eqref{eq:bound}, find the number of CG iterations
  guaranteed to reach relative energy-norm error $10^{-4}$ for $\kappa=4$,
  $\kappa=100$, and $\kappa=2500$. Confirm that each fourfold (then $25$-fold)
  increase in $\kappa$ roughly doubles (then triples) the count, matching the
  $\sqrt\kappa$ scaling.
  \item Write the CG step of \cref{sec:iterate} explicitly in the
  \code{iterate\_while\_with\_prev} form of \cref{ch:fold}: identify the bootstrap
  step (where $\vp_0=\vr_0$ and there is no previous $\beta$), the steady step, and
  the recurrence carry. Which stratum (\cref{ch:strata}) does the carry live in,
  and why is it wrong to store it in \code{OpParams}?
  \item Count the floating-point operations in one CG step as a function of $N$
  (count the apply abstractly as ``one matvec''). Then count one PCG step. By how
  much does preconditioning increase the per-step cost, and under what condition on
  the iteration-count reduction does PCG still win overall? Tie your answer to
  \cref{wex:count}.
  \item \emph{(Reduction order.)} Two machines run the same CG solve but sum the
  inner products in different orders. Argue that they may report different
  iteration counts yet both return iterates correct to the stated tolerance.
  Which fields of the \code{SolveResult} (\cref{ch:krylov}) are bit-reproducible
  across the two runs, and which are not? (Forward-reference
  \cref{app:numerics}.)
\end{enumerate}
